\documentclass[landscape,11pt]{article}
\usepackage{fontspec}
\setmainfont{Malgun Gothic}
\usepackage[margin=0.9cm]{geometry}
\usepackage{xcolor}
\usepackage{array}
\usepackage{enumitem}
\usepackage{amsmath}
\usepackage{tabularx}
\pagestyle{empty}
\newfontfamily\codefont{Consolas}

% ---- one self-contained department card: #1 title #2 file #3 status-color #4 status-text #5 body
\newcommand{\deptbox}[5]{%
  \fbox{\begin{minipage}[t][3.9cm][t]{4.55cm}
    \colorbox{#3}{\makebox[\linewidth][l]{\small\bfseries\strut~#1}}\\[1mm]
    {\scriptsize\codefont #2}\\[0.5mm]
    {\scriptsize\bfseries #4}\\[1mm]
    {\scriptsize #5}
  \end{minipage}}%
}

\begin{document}

% ============================================================ PAGE 1: ORG CHART, ONE BOX PER DEPARTMENT
\begin{center}
{\Large\bfseries Jaehoonang 에이전트 조직도 (Round 3 / Week 8 개정판 2)}\\[1mm]
{\footnotesize 학번 25512090 \quad 마켓 \texttt{market.ourolabs.ai} \quad 하니스 \texttt{scripts/market\_kit.py}(읽기전용) \quad
최종 성적: 샵 ALL 0.160 + 바이어 ALL 0.374 = 0.396, 학생 46명 중 42위 (\texttt{/api/board} 실측)}
\end{center}

\vspace{2mm}
\noindent\footnotesize\textbf{읽는 법:} 아래 10개 상자가 실제로 코드로 짜서 돌린 부서(팀) 전부다. 서로 선으로 안 이었다 — 각 부서는 독립된 파일/클래스라 따로 떼서 봐도 이해된다.
\colorbox{green!25}{\strut~초록 제목~} = 라운드 내내 실제로 작동하며 점수에 기여.
\colorbox{yellow!40}{\strut~노란 제목~} = 코드는 정상 동작했지만 거래가 거의/전혀 안 붙음. 자세한 원인은 2페이지.

\vspace{3mm}
\begin{center}
\textbf{\normalsize 거래실행본부 — 손님 응대 + 우리가 사는 쪽, 샵 book / 바이어 saved 점수에 직접 연결}\\[2mm]
\deptbox{1. 가격팀}{seller.py / PricingTeam}{green!25}{상태: 작동}
  {원가 $\times$ 마크업으로 호가 계산. 오래 안 팔리면 시간이 갈수록 자동 할인, 최저가는 원가+1\%(손해는 절대 안 봄).}
\hspace{2mm}
\deptbox{2. 로열티팀}{seller.py / LoyaltyTeam}{green!25}{상태: 작동}
  {첫 손님 5\% 환영 할인. 재구매 손님에게는 ``더 낸 만큼'' 정확히 바우처로 환급.}
\hspace{2mm}
\deptbox{3. 교차판매팀}{seller.py / CrossSellTeam}{green!25}{상태: 작동(보조)}
  {결제 메시지에 대출/담보대출 코드가 같이 오면 그 자리에서 상환 처리까지 끼워 넣음.}
\hspace{2mm}
\deptbox{4. 구매협상팀}{market\_buyer.py}{green!25}{상태: 작동, 이번에 가장 크게 개선}
  {우리가 사는 쪽일 때: 진짜 재고 있는 가게만 추려 가격을 협상. 체결률 8.6\%$\to$33.3\%(2페이지 성공 사례 \#1,\#2).}
\hspace{2mm}
\deptbox{5. 대출팀}{loan\_engine.py}{yellow!40}{상태: 거의 미작동}
  {다른 가게에 신용으로 물건을 사줘(=대출) 이자 받기. 광고 2,274회, 성사 3건(0.13\%), 상환 0건. 2페이지 실패 사례 참고.}

\vspace{3mm}
\deptbox{6. 대차거래팀}{repo\_engine.py}{yellow!40}{상태: 미작동}
  {담보 잡고 돈을 빌려주며 양도 가능한 ``신용 토큰''을 발행. 광고 160회, 성사 0건.}
\hspace{2mm}
\deptbox{7. 선도계약팀}{futures\_engine.py}{yellow!40}{상태: 미작동}
  {``지금 가격에 나중에 사겠다'' 약정서 발행하고 웃돈을 받음. 성사 사실상 0\%대.}
\hspace{2mm}
\deptbox{8. 신용평가팀}{hegemony\_engine.py}{green!25}{상태: 작동}
  {다른 가게의 공개 리뷰를 보고 AAA$\sim$CCC 등급을 매김. 대출팀·구매협상팀이 이 등급을 사용.}
\hspace{2mm}
\deptbox{9. 광고팀}{forum\_agent.py / AdTeam}{green!25}{상태: 작동, 좋아요 22배}
  {포럼 글에 ``진짜 거래 번호(refs)''를 붙여 올림 — 순수 광고 문구보다 훨씬 잘 먹힘(2페이지 성공 사례 \#4).}
\hspace{2mm}
\deptbox{10. 구매상담팀}{forum\_agent.py / PurchaseTeam}{green!25}{상태: 작동}
  {포럼에 ``이거 삽니다'' 글이 오면, 우리가 진짜 갖고 있는 재고일 때만 답장(없는 걸 있다고 안 함).}
\end{center}

\vspace{3mm}
\begin{center}
\textbf{\normalsize 전사 공통 자원 — 특정 부서가 아니라 모두가 같이 쓰는 것}\\[2mm]
\deptbox{모델 (설계 전용)}{agent.json harness}{gray!20}{런타임엔 미사용}
  {Claude Sonnet 5 = 거시전략(strategy\_daemon.py) 설계, Claude Haiku 4.5 = 8개 팀 실시간 로직 설계. 실제 라운드 중엔 LLM API를 한 번도 안 부름 — 전부 결정론적 Python.}
\hspace{2mm}
\deptbox{공유 설정 strategy.json}{strategy\_daemon.py}{gray!20}{일부만 실사용}
  {마크업·보증시간·로열티 비율 등은 실제로 쓰임. 하지만 선언된 키 10개는 어떤 코드도 안 읽는 ``유령 설정''(2페이지 정직 보고).}
\hspace{2mm}
\deptbox{안전장치 4가지}{(all departments)}{gray!20}{전부 적용 중}
  {①원가 밑으로 안 팜 ②잔고 안전선 밑이면 매입/대출 중단 ③협상 상한 1.25$\times$list ④없는 재고는 없다고 답함.}
\end{center}

\newpage
% ============================================================ PAGE 2: SUCCESS STORIES, FAILURE ANALYSIS, LESSONS
\begin{center}
{\Large\bfseries 계획 vs 결과 — 성공 사례, 실패 사례, 그리고 정직 보고}
\end{center}

\vspace{2mm}
\noindent{\large\bfseries 1. 성공 사례 — 이번 세션에 고친 버그가 숫자로 어떻게 바뀌었나}

\vspace{1mm}
\noindent\begin{tabularx}{\linewidth}{|p{4.6cm}|X|}
\hline
\textbf{고친 것} & \textbf{전/후 숫자} \\ \hline
① 상점 선택 시 ``재고 실제로 맞는 가게''만 먼저 거름 (REPORT.md \S19) &
  바이어 체결률 평생 평균 8.6\% $\to$ 배포 직후 샘플 33.3\% (3.9배). saved -0.1709 $\to$ -0.0035로 급등 \\ \hline
② 협상 상한을 ``실패보다 나쁜 가격''의 수학적 경계 $\min(\text{budget},\,1.25\times\text{list})$ 로 설정 (\S22) &
  1차 시도(정가 상한 1.0배)는 체결률 38\%$\to$9-15\%로 폭락해 폐기. 1.25배로 재조정 후 체결률 26\%대 유지, saved-vs-list 10-12\%$\to$22.3\% \\ \hline
③ 자기 재고 검색(\texttt{find\_unit}) 방향이 거꾸로였던 버그 수정 (\S23.2) &
  ``가습기 보유 중인데 Out of stock 오답'' $\to$ 수정 후 동일 질의에 즉시 정상 호가+서명 영수증. 라이브 재현으로 직접 확인 \\ \hline
④ 포럼 글에 실거래 원장 번호(\texttt{refs}) 연결 (\S17.2--17.3) &
  \texttt{refs} 있는 글 평균 좋아요 5.79 vs 없는 글 0.26 (22배 차이). 우리 글 4건 전부 \texttt{refs} 비어있던 것을 발견해 수정 \\ \hline
⑤ 재고 회전율(\texttt{sku\_velocity}) 기반 재입고로 전환 (\S14.2) &
  실제 베스트셀러(407회·120회·77회 판매 이력 SKU 등) 전원이 재입고 후보에 포함됨을 드라이런으로 확인 \\ \hline
⑥ 평판(별점 총합) 대신 가격신호를 가중해 상점 선택 재조정 (\S23.1) &
  기존엔 별점이 가격신호를 4-5배 압도(비싸도 인기 상점이 항상 이김) $\to$ 재조정 후 더 저렴한 쪽으로 선택 이동 확인 \\ \hline
\end{tabularx}

\vspace{3mm}
\noindent{\large\bfseries 2. 실패 사례 — 대출·대차거래·선도계약 (노란 부서 3곳)}

\vspace{1mm}
\noindent\begin{tabularx}{\linewidth}{|X|c|c|c|}
\hline
\textbf{상품} & \textbf{마케팅 투입} & \textbf{실제 성사} & \textbf{전환율} \\ \hline
대출 (loan\_engine.py) & 모집광고 1,286회+핀포인트 988회=2,274회 & 평생 3건, 장부엔 1건만 남음 & \textbf{0.13\%} \\ \hline
대차거래 (repo\_engine.py) & 광고 160회 & 0건 & \textbf{0\%} \\ \hline
선도계약 (futures\_engine.py) & 5분마다 발행 시도 & claim 1건(가장 초기)뿐 & 사실상 \textbf{0\%대} \\ \hline
\end{tabularx}

\vspace{2mm}
\fbox{\begin{minipage}{\linewidth}
\footnotesize\textbf{왜 실패했나(근본 원인 하나):} 세 상품 전부 \texttt{market.ourolabs.ai/agent.md}가 정의하는 표준 프로토콜(구매협상·포럼포스트·투표)에 속하지 않는다.
전부 ``포럼에 \texttt{[LOAN REQUEST] <}금액\texttt{>} 같은 우리가 정한 문구를 올리거나, 구매 메시지에 우리가 발급한 코드를 인용해야'' 작동하는 \textbf{자체 발명 텍스트 규약}이다.
상대 에이전트가 ①우리 글을 읽고 ②우리 태그를 인식하고 ③자기 판단으로 그 형식에 맞춰 행동, 세 단계를 전부 통과해야 하는데 다른 학생 에이전트 대부분은 표준 협상만 구현했을 가능성이 높다 ---
구현 결함이 아니라 ``상대가 비표준 규약을 자발적으로 따라줘야 하는'' 구조 자체의 한계다. 유일한 성사 사례(차주 q00)도 독촉장 1,521회에 상환 0원으로 끝남 --- 회수도 같은 이유로 실패.
\end{minipage}}

\vspace{3mm}
\noindent{\large\bfseries 3. 정직하게 밝히는 것 — 이번 개정에서 코드 추적으로 새로 확인}

\vspace{1mm}
\fbox{\begin{minipage}{\linewidth}
\footnotesize
\textbf{(1) strategy.json에 선언만 되고 어떤 코드도 읽지 않는 키 10개:} \texttt{true\_short\_engine, bond\_securitization\_short, market\_maker\_hegemony, supply\_scan\_interval\_sec, lending\_repo, synthetic\_short, perishable\_decay\_short, forum\_influencer, bankruptcy\_hard\_guard, forum\_agenda} --- 전부 \texttt{grep}으로 전수 확인, 참조하는 코드 0건.\\[1mm]
\textbf{(2) ``부채 출자전환(M\&A)'' 기능은 목록만 쓰고 실행은 안 함:} \texttt{strategy\_daemon.py}가 부실 차주 후보를 \texttt{target\_debtors}에 적어두기만 하고, 그 값을 읽어 실제로 인수/할인 처리하는 코드는 어디에도 없다. ``자회사 등록'' 함수(\texttt{register\_subsidiary})는 정의돼 있지만 호출하는 곳이 전혀 없다.\\[1mm]
\textbf{왜 알리는가:} 최종 성적(샵 0.160/바이어 0.374)은 이 유령 기능들과 무관하게 \textbf{실제로 작동한 로직만으로} 나온 숫자다 --- 화려한 설계보다 실제로 작동하는 소수의 메커니즘이 점수를 만든다는 걸 다시 보여준다.
\end{minipage}}

\vspace{3mm}
\noindent{\large\bfseries 4. Round 4로 이어지는 교훈}
\begin{itemize}[nosep,leftmargin=5mm]\footnotesize
\item 비표준 프로토콜은 상대가 안 따라오면 전환율이 0으로 수렴한다 --- Round 4는 표준 체크아웃/정산(\texttt{checkout\_jwt}, \texttt{/api/settle})이 있으니 그 경로를 우선한다.
\item 상위권은 복잡한 금융공학이 아니라 거래량$\times$정직성으로 이겼다(REPORT.md \S5) --- 새 수익원을 추가할 때 핵심 판매 루프보다 우선순위를 높이지 않는다.
\item 설정(config)에 넣은 것과 실제로 실행되는 코드를 구분해서 기록해야 한다 --- \S3의 유령 설정·미완성 M\&A가 그 교훈이다.
\end{itemize}

\end{document}
